\documentclass[11pt,oneside]{article}
\frenchspacing

\usepackage[utf8]{inputenc}
\usepackage[T1]{fontenc}
\usepackage{amsmath}
\usepackage{amssymb}
\usepackage{amsthm}
\usepackage{enumitem}
\usepackage{geometry}
\usepackage{hyperref}

\geometry{margin=1in}
\hypersetup{
  colorlinks=true,
  linkcolor=blue,
  urlcolor=blue,
}

\newtheorem{definition}{Definition}
\newtheorem{proposition}{Proposition}

\newcommand{\R}{\mathbb{R}}
\newcommand{\barp}{\bar p}
\newcommand{\estate}{e}

\title{Eisenberg--Noe Clearing as a Smart Contract Primitive}
\author{Research note}
\date{June 2026}

\begin{document}
\maketitle

\begin{abstract}
The research idea is to implement a self-clearing credit network on-chain.
Borrowers and lenders, or pools thereof in lending architectures such as
Morpho, Aave, and related venues, do not resolve distress through isolated
liquidations or first-come-first-served withdrawals. They opt into a common
liability graph.
When the network enters settlement, a smart contract applies an
Eisenberg--Noe style global bankruptcy rule: each debtor pays what it can, each
creditor receives pro rata within its claim class, and circular obligations are
cleared by a fixed point. The purpose of the contract is not to make default
disappear. It is to replace local liquidation races with a verifiable global
clearing vector.
\end{abstract}

\section{Motivation}

Most DeFi credit protocols are pairwise. A position becomes unsafe, a liquidator
buys collateral, and the local loan is closed. This is clean for a single
overcollateralized loan, but it is less natural for a network of vaults,
borrowers, lenders, rehypothecated claims, recursive lending loops, or tokenized
credit lines.

In such a network, a participant's solvency may depend on payments from other
participants who themselves depend on payments from the first participant or
from a shared set of intermediaries. Local liquidation can then destroy value
or amplify a stress that a global settlement would have absorbed.

The Eisenberg--Noe model gives a compact rule for this setting
\cite{eisenberg2001systemic}. It treats the financial system as a directed
graph of nominal obligations and computes the clearing payments that are
mutually consistent with limited liability and pro rata repayment. The
smart-contract question is:

\begin{quote}
Can we build an on-chain clearinghouse where the liability graph is native
state, the clearing vector is proposed by solvers, and the contract verifies
and executes the fixed point?
\end{quote}

\section{Debt Networks}

The older debt-network note starts from a useful observation: debt is a
directed graph, but solvency is not a radius-one property. A participant can
look solvent against its immediate creditors and debtors while still being
exposed to a payment failure farther away in the graph. Conversely, a
participant can look weak locally but survive if enough incoming payments clear.

This is the local/global tension. Contracts are written bilaterally, but
insolvency is settled globally. In a dense credit network, a lender does not
only need to know whether its direct borrower is solvent. It also needs to know
what hidden liabilities, pledged claims, and incoming payments support that
borrower. This is why rumours, reputation, disclosure, and regulation matter in
ordinary finance, and why observability and claim eligibility matter in DeFi.

Historical debt-entanglement cases make the same point. Repo market stress in
the 1980s showed that apparently safe repo chains can become system-wide
settlement problems once one participant fails \cite{garbade}. The Kuwait
al-Manakh stock-market crash generated a large post-crash debt-entanglement
problem that was treated by explicit clearing methods rather than isolated
bilateral enforcement \cite{elimam1997solution}.

For smart contracts, the debt-network view turns into three design questions:

\begin{itemize}[leftmargin=1.5em]
  \item What information about the graph must be public or committed before
  clearing?
  \item What is the objective of the clearing rule: value preservation, fast
  execution, fairness, minimum defaults, or something else?
  \item What liquidation or fire-sale externalities should the rule internalize
  before it starts selling assets into stress?
\end{itemize}

\section{The Clearing Model}

Let $N=\{1,\ldots,n\}$ be a set of participants. A debt network consists of
external assets $x_i\geq0$ and a raw liability matrix $L$. The entry
$L_{ij}\geq0$ is the amount participant $i$ has promised to pay participant
$j$, with $L_{ii}=0$. The total nominal payment due from $i$ is
\[
  \barp_i=\sum_j L_{ij}.
\]
The matrix $\Pi \in \R_+^{n \times n}$ records relative liabilities:
$\Pi_{ij}$ is the fraction of $i$'s total liabilities owed to $j$. If
$\barp_i>0$ then
\[
  \sum_j \Pi_{ij}=1.
\]
If $\barp_i=0$, set $\Pi_{ij}=0$ for all $j$.

A clearing payment vector $p \in \R_+^n$ satisfies
\[
  0 \leq p_i \leq \barp_i
\]
and
\[
  p_i =
  \min\left\{
    \barp_i,\,
    x_i + \sum_j \Pi_{ji}p_j
  \right\}.
\]
The term
\[
  \estate_i(p) = x_i + \sum_j \Pi_{ji}p_j
\]
is the estate available to participant $i$ after incoming payments. If
$\estate_i(p) \geq \barp_i$, participant $i$ pays all liabilities. If
$\estate_i(p)<\barp_i$, participant $i$ defaults and pays exactly its estate
pro rata to creditors.

The value left for participant $i$ after clearing is
\[
  v_i(p)=
  \left(x_i+\sum_j \Pi_{ji}p_j-\barp_i\right)_+.
\]
This makes clear what proportional clearing is doing. Creditors share the
payment of a defaulting node according to the liability shares $\Pi_{ij}$; any
surplus after full payment stays with equity. Symmetric bilateral netting
$L_{ij},L_{ji}\mapsto L_{ij}-L_{ji},0$ is therefore not always innocuous before
the clearing rule is applied, especially when one of the two nodes is near
default.

Equivalently, define the monotone map
\[
  \Phi_i(p)=
  \min\left\{
    \barp_i,\,
    x_i + \sum_j \Pi_{ji}p_j
  \right\}.
\]
A clearing vector is a fixed point $p=\Phi(p)$. Under standard regularity
conditions the relevant clearing vector is unique. Recent work sharpens the
equilibrium uniqueness question for the original model
\cite{stachurski2022properties}. Even when uniqueness needs care, the greatest
clearing vector has a clear economic meaning: it is the most optimistic payment
vector compatible with limited liability and the liability graph.

\section{Contract Interpretation}

The on-chain primitive is a settlement contract for a single numeraire asset at
first. Participants deposit assets, issue obligations to each other, and agree
that a settlement epoch will be resolved by the clearing rule.

\subsection{State}

The minimal state is:

\begin{itemize}[leftmargin=1.5em]
  \item a participant registry;
  \item a set of directed liability edges $(i,j,L_{ij})$;
  \item deposited external assets $x_i$ for each participant;
  \item a settlement status: live, frozen, proposed, challenged, or executed;
  \item the proposed clearing vector $p$ for the current epoch.
\end{itemize}

From the raw liability matrix $L$, the contract derives
\[
  \barp_i=\sum_j L_{ij}
  \quad\text{and}\quad
  \Pi_{ij} =
  \begin{cases}
    L_{ij}/\barp_i, & \barp_i>0,\\
    0, & \barp_i=0.
  \end{cases}
\]

\subsection{Lifecycle}

A simple lifecycle is:

\begin{enumerate}[leftmargin=1.5em]
  \item Live mode. Participants create, transfer, repay, or cancel obligations
  subject to protocol rules.
  \item Freeze. An epoch ends or a distress trigger freezes the graph.
  \item Proposal. A solver submits a clearing vector $p$ and posts a bond.
  \item Verification. The contract checks that $p$ satisfies the clearing
  equations within exact integer arithmetic or a stated tolerance.
  \item Challenge. Other participants can challenge an invalid vector during a
  short window.
  \item Execution. The contract distributes assets according to $p$ and marks
  unpaid liabilities as defaulted, rolled, tokenized, or extinguished according
  to the product design.
\end{enumerate}

\section{Why Verification Is Easier Than Solving}

The contract does not need to run an iterative fixed-point algorithm for large
graphs. A solver can compute $p$ off-chain. The contract only needs to verify
that the proposed vector is a fixed point.

For each participant $i$, compute
\[
  \estate_i = x_i + \sum_j L_{ji}\frac{p_j}{\barp_j},
\]
where the fraction is defined as zero if $\barp_j=0$. Then check:
\[
  0 \leq p_i \leq \barp_i.
\]
If $p_i < \barp_i$, require
\[
  p_i = \estate_i.
\]
If $p_i=\barp_i$, require
\[
  \estate_i \geq \barp_i.
\]
These local checks prove the global fixed-point condition. They cost
$O(n+|E|)$ for a graph with $n$ participants and $|E|$ liability edges.

\subsection{Certifying the greatest fixed point}

The canonical Eisenberg--Noe clearing vector is the greatest fixed point of
\[
  \Phi(p)=\barp \wedge (x+\Pi^\top p).
\]
By Tarski's theorem, the fixed points form a non-empty lattice. Economically,
the greatest fixed point is the maximal payment vector compatible with limited
liability, priority, and the liability graph. If the clearing vector is unique,
then any verified fixed point is already canonical. If uniqueness is not built
into the admissible graph, verifying
\[
  p=\Phi(p)
\]
only proves that $p$ is a fixed point. It does not by itself prove that $p$ is
the greatest one.

The direct certificate is the top-iteration trace:
\[
  p^0=\barp,\qquad p^{t+1}=\Phi(p^t).
\]
Since $\Phi$ is monotone and $\Phi(\barp)\leq \barp$, the sequence decreases.
In exact integer arithmetic with a deterministic rounding rule, the trace
terminates when
\[
  p^{T+1}=p^T.
\]
The terminal vector is the greatest fixed point. A solver can therefore submit
the whole trace
\[
  p^0,p^1,\ldots,p^T
\]
and the contract verifies every transition. This is stronger than a local
fixed-point check, but it costs $O(T|E|)$ rather than $O(|E|)$.

The same trace also defines default levels. A node is $0$-insolvent if it
cannot pay even when everyone else pays in full:
\[
  x_i+\sum_j L_{ji}<\barp_i.
\]
More generally, a node is $\mu$-insolvent if it first fails at the $\mu$-th
round of the top iteration. This gives a contagion clock. Nodes that default at
level $0$ are immediately short of resources; nodes that default later are
pulled into default by previous rounds of non-payment. For a soft-clearing
contract, these levels are natural warning states: margin calls can begin when
a node moves from solvent under $p^0$ to vulnerable under $p^1$ or $p^2$,
rather than waiting for the terminal fixed point.

There are several lighter variants:

\begin{itemize}[leftmargin=1.5em]
  \item enforce uniqueness conditions on admissible graphs, so a local
  fixed-point check is sufficient;
  \item use an optimistic challenge rule where any participant can disprove a
  proposal by submitting a strictly higher valid fixed point;
  \item require a proof that the submitted vector is the result of the
  top-iteration procedure, possibly with a succinct proof system;
  \item restrict the first implementation to graph classes where multiplicity is
  not the core risk.
\end{itemize}

This maximality issue may become less central in a soft-clearing design. If the
system moves through warning, compression, and partial deleveraging before hard
settlement, the protocol may care more about conservative executable recovery
values and monotone risk triggers than about proving a terminal bankruptcy
allocation in one step. Still, for a hard clearing epoch, the contract should
either verify the greatest fixed point or make explicit why any fixed point is
canonical.

\subsection{Integer Form}

In Solidity, the contract should avoid rational numbers. It can store raw
liabilities $L_{ij}$ and payments $p_i$ in integer token units. To check
participant $i$, multiply through by common denominators or use conservative
rounding rules.

For each incoming edge $j \to i$, the payment sent from $j$ to $i$ is
\[
  q_{ji} = \left\lfloor \frac{L_{ji}p_j}{\barp_j} \right\rfloor.
\]
The rounding remainder can be assigned to a deterministic residue account, sent
to creditors by edge order, or retained as dust. The important design point is
that the rounding rule must be part of the clearing equation, not an
implementation afterthought.

\section{Smart Contract Sketch}

The verifier has this shape:

\begin{verbatim}
for each participant i:
    require(p[i] <= totalLiabilities[i])
    estate = externalAssets[i]

for each liability edge j -> i with amount L[j][i]:
    if totalLiabilities[j] > 0:
        estate[i] += floor(L[j][i] * p[j] / totalLiabilities[j])

for each participant i:
    if p[i] < totalLiabilities[i]:
        require(p[i] == estate[i])
    else:
        require(estate[i] >= totalLiabilities[i])
\end{verbatim}

Execution then pushes edge payments:

\begin{verbatim}
for each liability edge i -> j with amount L[i][j]:
    if totalLiabilities[i] > 0:
        q = floor(L[i][j] * p[i] / totalLiabilities[i])
        transfer(asset, j, q)
\end{verbatim}

This is the core self-solving mechanism. Solvers compete to submit the correct
clearing vector. The contract enforces the bankruptcy rule. Participants no
longer need to race each other through local closeouts once the network is
frozen.

\section{What This Buys}

\subsection{Netting circular claims}

Circular obligations are not exceptional in the model. If $A$ owes $B$, $B$
owes $C$, and $C$ owes $A$, the fixed point clears the cycle jointly. The result
can preserve value that pairwise liquidation would waste.

\subsection{Contagion as state}

The default set is endogenous:
\[
  D(p)=\{i: p_i<\barp_i\}.
\]
Losses, contagion paths, and recovery rates become explicit outputs of the
settlement contract. This is useful for monitoring, insurance, tranche design,
and risk capital.

\subsection{Bankruptcy before liquidation}

The contract separates the insolvency rule from the asset liquidation rule.
In a collateralized DeFi market, collateral can be sold first and residual
unsecured claims can then enter the clearing graph. Or collateral can remain in
the estate and be distributed by the global clearing vector. These are different
products, but the same fixed-point logic can sit beneath both.

\section{DeFi Extensions}

The plain Eisenberg--Noe model assumes a single numeraire, fixed nominal
liabilities, and pro rata claims of equal seniority. DeFi needs extensions.
Close relatives in the financial-network literature add default costs, rescue
rules, liquidity effects, and network-stability analysis
\cite{rogers2013failure,cifuentes2005liquidity,acemoglu2015systemic,glasserman2016contagion}.

\subsection{Collateral}

Collateral can be handled by a waterfall:

\begin{enumerate}[leftmargin=1.5em]
  \item seize or mark collateral;
  \item repay secured creditors up to the collateral value;
  \item place any residual deficiency into the unsecured clearing graph;
  \item clear unsecured liabilities by the Eisenberg--Noe rule.
\end{enumerate}

An alternative is to treat collateral as part of $x_i$ and allow the global
clearing to distribute the estate. That is simpler but may conflict with
secured-creditor expectations.

\subsection{Default costs and fire sales}

Rogers--Veraart style extensions add default costs to the clearing process
\cite{rogers2013failure}. In reduced form, only a fraction of outside assets or
incoming payments survives when a node defaults. These coefficients are often
written as haircuts such as $\alpha,\beta\in[0,1]$.

For a smart contract this is not necessarily deadweight loss. Some part of the
haircut may be an auction spread, liquidation penalty, market-maker revenue, or
protocol-controlled reserve. The design question is where those costs go. If
they are burned by fire-sale execution, the clearing rule should try to avoid
them. If they are paid to a market-clearing mechanism that supplies liquidity in
stress, they become part of the incentive design.

\subsection{Seniority}

Multiple seniority classes can be implemented by sequential clearing. Senior
claims are cleared first. Remaining estate then enters the next class. This
preserves priority without abandoning the network view.

\subsection{Multiple assets}

A multi-asset version requires prices or exchange mechanisms. The cleanest MVP
uses a single ERC20 numeraire. A later version can introduce:

\begin{itemize}[leftmargin=1.5em]
  \item oracle prices fixed at the freeze block;
  \item on-chain auctions for non-numeraire collateral;
  \item separate clearing vectors by asset;
  \item endogenous price impact if liquidation size affects recovery.
\end{itemize}

\subsection{Repos and off-chain RWA claims}

Repo networks and tokenized RWA credit share a limitation that the plain model
does not capture: time and collateral are part of the claim. A repo is not just
an unsecured promise to pay. It is a dated funding transaction with collateral,
haircuts, closeout rights, and settlement mechanics. Likewise, an RWA token may
represent a claim on an issuer, SPV, receivable pool, custodian account, or fund
share.

The first on-chain implementation should therefore be explicit about what has
entered the clearing graph. A token balance can be a node in the graph, but the
off-chain estate behind that token must be represented by a separate attested
input: collateral documents, reserve reports, liquidation-value estimates,
seniority terms, and expected settlement delay. In that setting the clearing
contract clears the on-chain claim graph, while the off-chain asset graph
supplies estate values and waterfalls.

\subsection{Dynamic obligations}

The graph must freeze before settlement. Otherwise participants can create
claims after observing the proposed clearing vector. A real implementation needs
epoch boundaries, rate limits, permissioning, or a mempool-resistant freeze
rule.

\section{Objections as Design Constraints}

The natural objections narrow to two design questions:

\begin{itemize}[leftmargin=1.5em]
  \item which claims are eligible for the clearing graph;
  \item how secured and collateralized claims enter the clearing rule.
\end{itemize}

Some apparent objections are implementation choices rather than blockers.
Solving the fixed point on-chain is unnecessary: the solver can compute the
candidate vector off-chain, while the contract verifies the fixed-point
conditions in $O(n+|E|)$. The remaining subtlety is maximality: if uniqueness
conditions are not enforced by the admissible graph, the contract should require
the greatest clearing vector rather than merely any fixed point, using one of
the certificate or challenge rules above.

Oracle manipulation is also not special to this construction. Any DeFi system
that admits marked collateral inherits oracle risk. The sharper version is that
the oracle becomes a bankruptcy allocation input. The first prototype should
therefore use a single deposited numeraire and no oracle. Marked collateral can
enter only after the fixed-point verifier and settlement accounting are clean.

\subsection{Why plain pro rata bankruptcy is not enough}

Plain Eisenberg--Noe is pari passu unsecured bankruptcy. Each creditor of a
defaulting node shares the estate pro rata within a claim class. DeFi lending is
usually not shaped like that. It has secured claims, liquidation thresholds,
collateral-specific liens, seniority, liquidation penalties, and sometimes
price impact.

Consider a debtor $A$ with:

\[
\begin{array}{rcl}
  A & \text{owes} & B: 100, \quad \text{secured by collateral worth } 80,\\
  A & \text{owes} & C: 100, \quad \text{unsecured},\\
  A & \text{has} & 20 \quad \text{free assets}.
\end{array}
\]

Plain pro rata clearing would see estate $100$ and liabilities $200$, so $B$
receives $50$ and $C$ receives $50$. That ignores the secured nature of $B$'s
claim.

A secured-credit waterfall gives a different result. First, $B$ receives the
collateral value $80$. This leaves a residual unsecured deficiency of $20$ for
$B$. The unsecured graph then has claims $B:20$ and $C:100$ against free estate
$20$, so $B$ receives $20\cdot 20/120=3.33$ and $C$ receives
$20\cdot 100/120=16.67$. The final allocation is:

\[
  B:83.33,\qquad C:16.67.
\]

This suggests the clean extension:

\begin{enumerate}[leftmargin=1.5em]
  \item resolve secured closeout or collateral waterfalls first;
  \item convert collateral shortfalls into unsecured deficiency claims;
  \item add any collateral surplus back to the debtor estate;
  \item run Eisenberg--Noe on the residual unsecured graph.
\end{enumerate}

The fixed-point logic is then preserved, but it applies to the residual claims
that should actually be pari passu.

\subsection{Strategic claim creation near clearing}

The strongest strategic issue is not hidden bilateral obligations. Hidden claims
should simply be ineligible for the current clearing. The dangerous case is a
visible but late obligation or claim transfer inserted once participants know
that clearing is near.

Examples include:

\begin{itemize}[leftmargin=1.5em]
  \item a distressed borrower creates a friendly liability to an affiliated
  address;
  \item a creditor transfers claims to a party with different treatment;
  \item claims are routed through an entity expected to default in order to
  change the loss allocation;
  \item a participant front-runs the freeze trigger with a new edge.
\end{itemize}

The defense is an eligibility and seasoning rule:

\begin{quote}
Only obligations created before snapshot block $B-k$ enter the current
clearing. Late obligations roll into the next epoch. Claim transfers preserve
the original claim age and seniority. New obligations must arise from
protocol-verified value transfer, not arbitrary IOU creation.
\end{quote}

This makes the clearing graph a state machine rather than a free-form list of
asserted debts.

\subsection{Fission, fusion, and strategic manipulation}

The old debt-network note points to a separate mechanism-design problem: nodes
may want to split or merge before clearing. Let $f:N\to M$ be a map that groups
old nodes into new nodes. A fusion operator sends a debt network $(x,L)$ on
$N$ to a coarser network on $M$ by
\[
  \widehat f(x)_a=\sum_{i\in f^{-1}(a)}x_i
\]
and, for $a\neq b$,
\[
  \widehat f(L)_{ab}
  =
  \sum_{i\in f^{-1}(a)}\sum_{j\in f^{-1}(b)} L_{ij}.
\]
Fusion is the mathematical version of a rescue consortium: several nodes pool
assets and liabilities so that a quasi-solvent group can support an insolvent
member. This can preserve value when liquidation costs are high.

Fission is the adversarial direction. If a distressed borrower can split into
several addresses, assign liabilities selectively, and create a large friendly
claim to an affiliated sink, proportional repayment can redirect estate away
from outside creditors. This is a Sybil problem for bankruptcy rules. The
contract therefore needs identity, seasoning, and provenance constraints:
claim transfers should preserve original age and seniority, internal claims
should be marked, and arbitrary IOUs should not become eligible claims merely
because they were written to the graph before freeze.

This is a deeper version of the late-claim problem. It asks whether the clearing
mechanism is fusion-proof and fission-proof, or whether it at least makes the
attack surface explicit enough to control through registry rules.

A few useful facts guide the mechanism design. With zero fusion cost, fusing two
solvent nodes preserves solvency: the fused node pools assets, incoming claims,
and outgoing liabilities, and internal debts disappear. With positive fusion
cost $\kappa$, the same statement only holds if the pooled surplus exceeds
$\kappa$. At the other extreme, fusing two nodes that are already
$0$-insolvent does not by itself restore solvency; their combined immediate
estate is still below their combined external liabilities after internal debts
are netted. More generally, fusing two nodes with the same default level cannot
be assumed to stop contagion unless the fusion brings in enough surplus or
removes enough cost.

The interesting case is mixed-level fusion. A quasi-solvent node may want to
fuse with, lend to, or rescue an earlier-defaulting node because the rescue
preserves incoming payments elsewhere in the graph. This can be socially useful:
it is the algorithmic analogue of a rescue consortium. But it can also be a
side-payment game that changes creditor priority. The rule should therefore
distinguish three cases:

\begin{itemize}[leftmargin=1.5em]
  \item admissible fusion, where the combined entity assumes all external
  liabilities and improves or preserves every external creditor's recovery;
  \item rescue lending, where a new value transfer enters the debtor estate but
  the lender does not jump the creditor queue;
  \item collusive fission or fake debt, where internal claims are manufactured
  to redirect an insolvent estate away from outside creditors.
\end{itemize}

\subsection{Soft deleveraging before hard clearing}

The network does not need to jump directly from normal lending to bankruptcy
settlement. Clearing can be soft and pre-emptive:

\begin{center}
\begin{tabular}{ll}
normal mode & ordinary borrowing, lending, repayment, and transfers,\\
warning mode & risk limits tighten and new concentration is discouraged,\\
compression mode & new debt is disabled and repayments remain open,\\
freeze mode & the eligible graph is snapshotted,\\
clearing mode & the clearing vector is verified and executed.
\end{tabular}
\end{center}

Triggers can be scheduled or risk-based: projected default under a haircut,
network leverage above a bound, collateral liquidity below a bound, oracle
deviation, failed margin top-up, participant petition, or governance call.

\subsection{Dynamic graphs and self-stabilisation}

The old note's ``self-stabilising debt network'' idea is the continuous version
of soft clearing. Instead of waiting for a terminal bankruptcy epoch, the
contract recomputes recovery values or clearing stress indicators over time and
issues margin calls before the graph becomes unsolvable.

The relevant events are:

\begin{itemize}[leftmargin=1.5em]
  \item node birth and death;
  \item creation, transfer, ageing, and deletion of liability edges;
  \item repayment and cash-flow events;
  \item market events that change collateral or liquidation value;
  \item compression, margin call, and clearing events.
\end{itemize}

This dynamic view is the bridge between Eisenberg--Noe as bankruptcy and
Eisenberg--Noe as risk control. The clearing vector becomes a continuously
computed stress state, not only a terminal payout.

\subsection{A candidate clearing rule}

A concrete first rule is:

\begin{enumerate}[leftmargin=1.5em]
  \item Snapshot the liability graph and admitted collateral at block $B$.
  \item Exclude non-seasoned claims from the current epoch.
  \item Close secured claims against their specific collateral first.
  \item Convert collateral shortfalls into unsecured deficiency claims.
  \item Add collateral surplus back to the debtor estate.
  \item Run Eisenberg--Noe on the residual unsecured graph.
  \item Execute pro rata payments and, if useful, tokenize unpaid residual
  claims.
\end{enumerate}

The rule keeps collateral priority explicit while still using global clearing
for the part of the system where global clearing is the correct object:
unsecured residual claims and circular payment dependencies.

\section{Research Questions}

\begin{itemize}[leftmargin=1.5em]
  \item What is the right trigger: scheduled epoch, solvency breach, governance
  call, oracle event, or participant petition?
  \item How should solver bonds and challenges be designed so invalid vectors
  are cheap to reject and valid vectors are cheap to finalize?
  \item Can the verifier be compressed with Merkle commitments, fraud proofs, or
  zero-knowledge proofs for large graphs?
  \item How do we handle collateral priority without reducing the system back to
  isolated liquidations?
  \item Which DeFi liabilities are eligible: vault shares, debt tokens, margin
  accounts, credit lines, receivables, or RWA notes?
  \item What strategic behavior appears when agents know that obligations will
  be globally cleared?
  \item Is the clearing rule robust to fission and fusion of nodes, or does it
  need a registry that fixes economic identity before settlement?
  \item Can useful clearing be computed on a visible subgraph, or does the rule
  require a panoptic view of the whole debt network?
  \item Can privacy, commitments, or peer-surveillance reveal enough of the
  graph without publishing every bilateral exposure?
  \item Can streams and rescheduling distinguish illiquidity from insolvency in
  a way that a single maturity payment cannot?
  \item Can the clearing outputs feed insurance premia, capital buffers, or
  curator risk limits?
\end{itemize}

\section{Minimal Viable Experiment}

The first useful implementation should be deliberately narrow:

\begin{itemize}[leftmargin=1.5em]
  \item single ERC20 numeraire;
  \item no oracle;
  \item all external assets are deposited in the clearing contract;
  \item fixed participant set during each epoch;
  \item raw liability edges registered on-chain;
  \item off-chain solver;
  \item on-chain fixed-point verification;
  \item optimistic challenge window;
  \item deterministic integer rounding.
\end{itemize}

This version is enough to test the central claim: a smart contract can enforce a
global bankruptcy rule over a liability graph without solving the fixed point
on-chain.

\section{Connection to Recursive Lending}

The model is especially relevant to recursive vault and lending structures.
When a vault lends against its own shares, or when vault shares become collateral
inside another lending venue, the dependency graph can contain cycles. Pairwise
liquidation sees only the local edge. A global clearing contract sees the cycle.

This suggests a product direction:

\begin{quote}
Build a credit network where recursive DeFi positions opt into a common
settlement layer. Normal times use ordinary lending mechanics. Stress times use
global clearing.
\end{quote}

The resulting system would sit between two familiar designs. It is not a simple
AMM or isolated lending market. It is also not a discretionary bankruptcy court.
It is a rule-based clearing court implemented as a contract.

\section{Pre-Clearing Games and Standstill Rules}

The Theotime note reframes clearing as a game played before the clearing block.
Players may still be able to lend, repay, withdraw assets, or create bilateral
claims after they understand that clearing is likely. These moves preserve or
redistribute total assets in different ways, and the clearing vector then turns
them into payoffs.

Two elementary moves are enough to expose the issue. If $i$ lends $x$ to $j$,
then
\[
  e_i\leftarrow e_i-x,\qquad e_j\leftarrow e_j+x,\qquad
  L_{ji}\leftarrow L_{ji}+x.
\]
If $j$ repays $y$ to $i$, then
\[
  e_j\leftarrow e_j-y,\qquad e_i\leftarrow e_i+y,\qquad
  L_{ji}\leftarrow L_{ji}-y.
\]
Total assets are unchanged, but the liability graph and the later proportional
allocation are changed. This is why a pre-clearing standstill matters.

The danger is not only fraud. Some pre-clearing moves can be Pareto-improving
relative to an immediate hard freeze. Early repayment can reduce default costs
or fire-sale losses. In a simple default-cost model, if an insolvent node $0$
uses $x_i$ to repay creditor $i$ before clearing, system value can rise because
less value passes through the default haircut. In reduced form, the marginal
effect on aggregate recovery has the sign of
\[
  1-\alpha,
\]
where $\alpha$ is the outside-asset recovery coefficient. If $\alpha<1$, a
repayment before the default haircut can preserve value for the system.

The distributional effect is different. A prepayment to one creditor can hurt
other creditors by violating proportionality. This creates the central
standstill tradeoff:

\begin{itemize}[leftmargin=1.5em]
  \item a total freeze protects the proportionality principle but may destroy
  value by forcing default-cost or fire-sale losses;
  \item a permissive pre-clearing game may preserve value but lets insiders
  bargain around the global priority rule;
  \item a controlled standstill should allow value-preserving repayments or
  rescue funding only if they are public, seasoned, and do not lower any
  external creditor's recovery.
\end{itemize}

This suggests a contract rule: after warning mode begins, ordinary repayments
remain allowed, but preferential repayments, new insider claims, and bilateral
side payments require either creditor consent, a no-harm proof against the
current stress vector, or automatic conversion into junior rescue claims. This
is the mechanism-design version of soft clearing.

\section{Toy Attacks and Rescue Games}

The feed notes contain several small games that should eventually become
worked examples. The point of the examples is not to replace the clearing rule.
It is to show why the rules around claim creation, rescue funding, and
standstill matter.

\subsection{Default-cost multiplicity}

With default costs, the greatest fixed point matters even more. In a
Rogers--Veraart style extension, a defaulting node pays only a fraction
$\alpha$ of outside assets and a fraction $\beta$ of incoming payments. Consider
two nodes with symmetric liabilities
\[
  L=
  \begin{pmatrix}
    0 & \gamma\\
    \gamma & 0
  \end{pmatrix},
  \qquad
  x=
  \begin{pmatrix}
    1\\
    1
  \end{pmatrix},
  \qquad
  \alpha=\beta=\frac12,
  \qquad
  \gamma=2.2.
\]
Then both
\[
  p=
  \begin{pmatrix}
    2.2\\
    2.2
  \end{pmatrix}
  \quad\text{and}\quad
  p=
  \begin{pmatrix}
    1\\
    1
  \end{pmatrix}
\]
can satisfy the clearing equations of the default-cost model. The upper vector
is the economically relevant one: it Pareto-dominates the lower vector. This is
the clean toy reason why a contract cannot merely accept any fixed point once
default costs or rescue rules enter the model.

\subsection{Fake-debt sink}

Take three nodes. Node $1$ is distressed. It owes $L_{12}$ to node $2$ and
$L_{1n}$ to node $n$. Node $2$ owes $L_{21}$ to node $1$. If node $n$ can create
a new eligible claim $S_1$ against node $1$ just before clearing, then node
$n$'s share of node $1$'s default estate becomes
\[
  \frac{L_{1n}+S_1}{L_{1n}+S_1+L_{12}}.
\]
This expression is increasing in $S_1$. In the limit, the new claimant can
redirect most of the estate that would otherwise have been shared with node
$2$. With a default-cost coefficient $\beta$, the diverted amount can approach
a term of order $\beta L_{21}$ in the simplest versions of the example.

This is not value creation. It is priority theft unless $S_1$ came with genuine
new rescue value and the new claim is treated consistently with the creditor
waterfall. The contract implication is direct: late insider claims should not
enter the same class as seasoned external claims.

\subsection{Early repayment}

Now take one insolvent node $0$ and solvent creditors $1,\ldots,n$. Suppose
node $0$ makes early repayments $x_i$ before hard clearing, with
\[
  \sum_i x_i\leq e_0,\qquad x_i\leq L_{0i}.
\]
If the remaining estate of node $0$ is later hit by default costs, then early
repayment can preserve aggregate value. In a reduced form calculation, the
aggregate recovery has marginal effect
\[
  \frac{\partial}{\partial x_i}\sum_k R_k = 1-\alpha.
\]
So when $\alpha<1$, early repayment can be system-improving.

But the individual recovery $R_i$ need not move in the same direction for every
creditor. Early repayment to one creditor can lower another creditor's recovery
by shrinking the estate that remains to be shared proportionally. This is the
standstill problem in its smallest form: the move can be good for total value
and bad for the priority rule.

\subsection{Rule suggested by the toys}

The toy games suggest a three-way classification:

\begin{itemize}[leftmargin=1.5em]
  \item seasoned obligations enter the current clearing class;
  \item genuine rescue funding is allowed, but it enters a separately specified
  rescue class rather than silently diluting old creditors;
  \item late fake debt, insider preference, and unseasoned claim transfers are
  either excluded from the epoch or subordinated.
\end{itemize}

That classification is a better primitive than a blanket freeze. It preserves
room for Pareto-improving rescue while blocking the easiest pre-clearing
manipulations.

\section{Temporal Debt Streams and Rescheduling}

The plain model treats each liability as a single nominal amount due at
clearing. The old note suggests a richer object: a debt can be a finite stream
of promised payments. Mathematically, the first finite-stream model can live in
$c_{00}$, the space of finitely supported sequences, with componentwise order.
More general sequence spaces such as $c_0$ or $\ell_\infty$ raise lattice and
closure questions \cite[chapter 16]{aliprantis}, but the MVP should start with
finite streams.

With streams, clearing is no longer only a question of who pays how much. It is
also a question of when payments are due. This makes rescheduling a native
operation. A node may be insolvent if every claim is accelerated today, but
solvent if cash flows can be flattened over time at an agreed discount curve.
That distinction is central for off-chain credit, repos, and RWA collateral,
where the legal claim may be good but slow.

One can define a valuation functional
\[
  \operatorname{val}(x)=\sum_t \gamma^t x_t
\]
for a fixed discount factor $\gamma$, then compare a compressed payment with a
longer repayment schedule. This turns hard clearing into a family of possible
workouts:

\begin{itemize}[leftmargin=1.5em]
  \item immediate closeout;
  \item maturity extension;
  \item coupon reduction;
  \item recovery-token issuance;
  \item dynamic margin calls on the stream value.
\end{itemize}

This extension should stay separate from the first contract, but it is likely
the right language for a system that wants to clear RWA claims or repo-like
positions without forcing all assets into same-block liquidation.

\section{Near-Term To-Dos}

\begin{itemize}[leftmargin=1.5em]
  \item Write the single-numeraire contract invariant precisely.
  \item Work through a three-node cycle example by hand.
  \item Specify the integer rounding rule and dust ownership.
  \item Decide whether the first prototype should execute payments directly or
  mint recovery tokens.
  \item Compare optimistic verification, full on-chain verification, and zk
  verification.
  \item Connect the graph input to YTD or DIG-style dependency extraction.
  \item Specify identity, seasoning, and internal-claim rules against fission
  attacks.
  \item Decide how liquidative values are computed and challenged.
\end{itemize}

\section{Varia}

For unwinding thresholds, the relevant asset value is not necessarily
mark-to-market value. A clearing network should consider computing
liquidative values: what each position can actually realize when it must be
unwound through the available exit routes.

This matters because a position can look solvent at quoted prices while being
insolvent at executable exit depth. The clearing trigger should therefore be
able to use haircut or liquidation-value estates,
\[
  x_i^{\mathrm{liq}} \leq x_i^{\mathrm{mtm}},
\]
where $x_i^{\mathrm{liq}}$ reflects pool depth, redemption queues, auction
discounts, oracle delay, venue concentration, and expected slippage under the
size of the unwind. In that version, the warning and compression modes are not
triggered by accounting insolvency alone. They are triggered when the network
approaches the point where executable liquidation value no longer covers the
eligible claims.

\bibliographystyle{plain}
\bibliography{Eisenberg-Noe}

\end{document}
